\subsection{Universally measurable functional calculus}

In this section the Hilbert spaces considered are complex.

Let u be a partial normal operator on a Hilbert space E. We denote by \(\mu_{x,y}\) (resp. \(\mu_{y}\) ) the spectral measure of \((x,y)\in\mathrm{E}\times\mathrm{E}\) (resp. of y) relative to u.

Let \(y \in \mathrm{E}\) . The map \(\mathcal{K}(\mathrm{Sp}(u)) \to \mathrm{E}\) defined by \(f \mapsto f(u)(y)\) satisfies

\[
\| f (u) (y) \| ^ {2} = \int_ {\mathrm{Sp} (u)} | f | ^ {2} \mu_ {y} = \| f \| _ {\mathcal {L} ^ {2} (\mathrm{Sp} (u), \mu_ {y})} ^ {2}.
\]

There therefore exists a unique isometric linear map \(ev_{y}\) from the space \(\mathrm{L}^{2}(\mathrm{Sp}(u),\mu_{y})\) into E such that \(\mathrm{ev}_{y}(\widetilde{f})=f(u)(y)\) if \(\widetilde{f}\) is the class of a function \(f\in\mathcal{K}(\mathrm{Sp}(u))\) .

Let \(f\) be a universally measurable function defined on the spectrum of \(u\) . We denote by \(\mathrm{D}_{f}\) the set of elements \(y \in \mathrm{E}\) such that \(f\) belongs to \(\mathcal{L}^{2}(\mathrm{Sp}(u), \mu_{y})\) .

PROPOSITION 3. --- Let f be a universally measurable function defined on the spectrum of u. The set D \(_{f}\) is a dense vector subspace of E. The map y ↦ ev \(_{y}\) (f) is a partial normal operator on E whose domain is D \(_{f}\) , and which is denoted f(u).

For every \(x \in \mathrm{E}\) and every \(y \in \mathrm{D}_f\) , we have \(f \in \mathcal{L}^1(\mathrm{Sp}(u), \mu_{x,y})\) and

\[
\langle x \mid f (u) y \rangle = \int_ {\operatorname{Sp} (u)} f \mu_ {x, y}.\tag{6}
\]

We shall prove the following more precise statement:

PROPOSITION 4. --- Let f be a universally measurable function on the spectrum of u.

\begin{enumerate}
\def\labelenumi{\alph{enumi})}
\item
  The set \(D_{f}\) is a dense vector subspace of E. The map \(f(u): y \mapsto \operatorname{ev}_{y}(f)\) from \(D_{f}\) into E is linear and coincides with \(1_{E}\) if f = 1;
\item
  Let \(g \in \mathcal{L}_{\mathrm{u}}(\mathrm{Sp}(u))\) . For every \(y \in \mathrm{D}_f\) and every \(x \in \mathrm{D}_g\) , we have \(fg \in \mathcal{L}^1(\mathrm{Sp}(u), \mu_{x,y})\) , and
\end{enumerate}

\[
\langle g (u) x \mid f (u) y \rangle = \int_ {\mathrm{Sp} (u)} \overline {{g}} f \mu_ {x, y}.\tag{7}
\]

\begin{enumerate}
\def\labelenumi{\alph{enumi})}
\setcounter{enumi}{2}
\tightlist
\item
  Suppose that \(E = L^{2}(X, \mu)\) where X is a locally compact topological space and \(\mu\) is a positive measure on X, and that \(u = m_{h}\) , where \(h: X \to \mathbf C\) is \(\mu\) -measurable. We have \(f(m_{h}) = m_{f \circ h}\) .
\end{enumerate}

By Theorem 1 of IV, p. 266, we may assume that we are in the situation of assertion \(c\) ), namely, \(\mathrm{E} = \mathrm{L}^2 (\mathrm{X},\mu)\) and \(u = m_h\) , where X is a locally compact topological space, \(\mu\) is a positive measure on X and \(h\colon \mathrm{X}\to \mathbf{C}\) is \(\mu\) -measurable. We shall denote by \(\widetilde{\varphi}\) the class in \(\mathrm{L}^2 (\mathrm{X},\mu)\) of a function \(\varphi \in \mathcal{L}^2 (\mathrm{X},\mu)\) .

Let \(\mathrm{S} = \mathrm{Sp}(m_h)\) . We denote by \(\mu_{\widetilde{\varphi}_1, \widetilde{\varphi}_2}\) (resp. \(\mu_{\widetilde{\varphi}}\) ) the spectral measure of \((\widetilde{\varphi}_1, \widetilde{\varphi}_2)\) (resp. of \(\widetilde{\varphi}\) ) relative to \(m_h\) for all \((\varphi_1, \varphi_2) \in \mathcal{L}^2(\mathrm{X}, \mu)^2\) (resp. every \(\varphi \in \mathcal{L}^2(\mathrm{X}, \mu)\) ).

Let \(\varphi\in\mathcal{L}^{2}(X,\mu)\) . The spectral measure \(\mu_{\widetilde{\varphi}}\) is equal to the image measure \(h(|\varphi|^{2}\cdot\mu)\) on S (Lemma 4 of IV, p. 269). Since the function \(f\) is \(\mu_{\varphi}\) -measurable, we have \(\varphi\in D_{f}\) if and only if the integrals

\[
\int_ {\mathrm{S}} ^ {*} | f | ^ {2} d \mu_ {\widetilde {\varphi}} = \int_ {\mathrm{S}} ^ {*} | f | ^ {2} h (| \varphi | ^ {2} d \mu) = \int_ {\mathrm{X}} ^ {*} | f \circ h | ^ {2} | \varphi | ^ {2} d \mu
\]

are finite (INT, V, p. 70, § 6, no. 2, prop. 2). This means that \(D_{f}\) is the domain of the multiplication operator \(m_{f\circ h}\) , which is a dense subspace of E (prop. 5 of IV, p. 232).

The restriction of \(\mathrm{ev}_{\widetilde{\varphi}}\) to the function classes \(g \in \mathcal{K}(\mathrm{S})\) is the map which associates to a class \(\widetilde{g}\) the element \(g(m_h)(\widetilde{\varphi}) = m_{g \circ h}(\widetilde{\varphi})\) (Lemma 4 of IV, p. 269). The map \(\mathrm{ev}_{\widetilde{\varphi}}\) therefore coincides with the isometric map from \(\mathrm{L}^2(\mathrm{S}, \mu_{\widetilde{\varphi}})\) into \(\mathrm{L}^2(\mathrm{X}, \mu)\) obtained by passing to quotients from the map \(g \mapsto (g \circ h) \cdot \varphi\) from \(\mathcal{L}^2(\mathrm{S}, \mu_{\widetilde{\varphi}})\) into \(\mathcal{L}^2(\mathrm{X}, \mu)\) .

In particular, we therefore have \(\mathrm{ev}_{\widetilde{\varphi}}(f)=m_{f\circ h}(\widetilde{\varphi})\) , so the map \(f(m_{h})\) from \(D_{f}\) into E coincides with the partial operator \(m_{f\circ h}\) . This proves assertion c); if f=1, we find \(\mathrm{ev}_{\widetilde{\varphi}}(1)=\widetilde{\varphi}\) , which also completes the proof of a).

Let us prove assertion \(b\) ). Let \(\varphi_{1}\) and \(\varphi_{2}\) be functions such that \(\widetilde{\varphi}_{1} \in \mathrm{D}_{f}\) and \(\widetilde{\varphi}_{2} \in \mathrm{D}_{g}\) . It follows from (INT, V, loc. cit.)

\[
\begin{array}{l} \int_ {\mathrm{S}} ^ {*} | f g |   | \mu_ {\widetilde {\varphi_ {1}}, \widetilde {\varphi_ {2}}} | = \int_ {\mathrm{X}} ^ {*} | (f \circ h) (g \circ h) \varphi_ {1} \varphi_ {2} | d \mu \\ \qquad \leqslant \| (f \circ h) \varphi_ {1} \| \| (g \circ h) \varphi_ {2} \| \end{array}
\]

which is finite, hence \(fg \in \mathcal{L}^1(\mathrm{S}, \mu_{\widetilde{\varphi}_1, \widetilde{\varphi}_2})\) . Moreover, we then have

\[
\begin{array}{c} \langle g (m _ {h}) (\widetilde {\varphi} _ {2})   |   f (m _ {h}) (\widetilde {\varphi} _ {1}) \rangle = \int_ {\mathrm{X}} \overline {{(g \circ h)   \varphi_ {2}}}    (f \circ h)   \varphi_ {1}    d \mu \\ = \int_ {\mathrm{S}} \overline {{g}}   f    d \mu_ {\widetilde {\varphi} _ {1}, \widetilde {\varphi} _ {2}}, \end{array}
\]

which concludes the proof.

DEFINITION 5. --- The map \(f \mapsto f(u)\) from \(\mathcal{L}_{\mathrm{u}}(\mathrm{Sp}(u))\) into the set of normal partial operators on E defined by prop. 3 of IV, p. 271 is called the universally measurable functional calculus mapping associated with \(u\) .

More generally, let T be a set and let \(g\colon\operatorname{Sp}(u)\times\operatorname{T}\to\mathbf{C}\) be a map; let \(t\in T\) be such that the function \(g_{t}\colon z\mapsto g(z,t)\) is universally measurable on \(\operatorname{Sp}(u)\) . We then denote by \(g(u,t)=g_{t}(u)\) .

Formula (7) implies in particular that

\[
\| f (u) (y) \| ^ {2} = \int_ {\operatorname{Sp} (u)} | f | ^ {2} \mu_ {y}\tag{8}
\]

for all \(f \in \mathcal{L}_{\mathrm{u}}(\mathrm{Sp}(u))\) and every \(y \in \mathrm{D}_f\) . From this, taking \(f = 1\) we deduce that \(\mu_y\) is a positive measure of total mass \(\| y\|^2\) .

If \(u\) is a normal partial operator on a Hilbert space E and \(f \in \mathcal{K}(\mathrm{Sp}(u))\) , then \(\mathrm{D}_f = \mathrm{E}\) and the partial operator \(f(u)\) is then continuous, since

\[
\| f (u) y \| \leqslant \| f \| _ {\infty} \| y \|
\]

for all \(y \in E\) according to (8). The endomorphism \(f(u)\) coincides with that of Definition 3 of IV, p. 268.

COROLLARY 1. --- a) For every \(f \in \mathcal{L}_{\mathrm{u}}(\mathrm{Sp}(u))\) , the partial operator \(f(u)\) is normal and \(f(u)^{*} = \overline{f}(u)\) . Moreover \(f(u)\) is positive if \(f \geqslant 0\) , and self-adjoint if \(f\) takes real values;

\begin{enumerate}
\def\labelenumi{\alph{enumi})}
\setcounter{enumi}{1}
\item
  Let \(k \in \mathbf{N}\) and \(f(z) = z^{k}\) for \(z \in \operatorname{Sp}(u)\) . We have \(f(u) = u^{k}\) ;
\item
  Let \(\lambda \in \mathbf{C} - \mathrm{Sp}(u)\) and \(f(z) = (\lambda - z)^{-1}\) for \(z \in \mathrm{Sp}(u)\) . We have \(f(u) = \mathrm{R}(u, \lambda)\) ;
\item
  One has \(\beta(u) = b(u)\) ;
\item
  Let \(f \in \mathcal{L}_{\mathrm{u}}(\mathrm{Sp}(u))\) and \(g \in \mathcal{L}_{\mathrm{u}}(\mathrm{Sp}(u))\) such that \(|g| \leqslant 1 + |f|\) . The domain of \(g(u)\) is a core for \(f(u)\) .
\end{enumerate}

In view of the spectral theorem (th. 1 of IV, p. 266), this follows from the preceding proposition combined, respectively, with:

\begin{enumerate}
\def\labelenumi{\alph{enumi})}
\item
  lemma 2, a) of IV, p. 266, prop. 23 of IV, p. 253, and its corollary;
\item
  prop. 24, b) of IV, p. 255;
\item
  Prop. 22, b) of IV, p. 252;
\item
  Lemma 2, a) of IV, p. 266;
\item
  Prop. 6, b) of IV, p. 232.
\end{enumerate}

Note. --- For \(f = \operatorname{Id}_{\operatorname{Sp}(u)}\) , we have \(f(u) = u\) (assertion b) for \(k = 1\) ). The domain of u therefore coincides with the set of \(x \in E\) such that the identity map of \(\operatorname{Sp}(u)\) belongs to \(\mathcal{L}^{2}(\operatorname{Sp}(u), \mu_{x})\) ; it contains in particular the elements \(x \in E\) such that the measure \(\mu_{x}\) is compactly supported.

The following corollary generalizes the corollary of Prop. 16 of IV, p. 247 and Prop. 17 of IV, p. 248.

COROLLARY 2. --- Let u be a partial normal operator on the Hilbert space E. Assume that E is nonzero.

\begin{enumerate}
\def\labelenumi{\alph{enumi})}
\item
  The spectrum of u is nonempty;
\item
  For every \(\lambda \in \mathbf{C} - \mathrm{Sp}(u)\) , the norm of the resolvent \(\mathrm{R}(u, \lambda)\) is equal to \(1/\delta\) , where \(\delta > 0\) is the distance in \(\mathbf{C}\) from \(\lambda\) to the spectrum of \(u\) .
\item
  For every \(\varepsilon > 0\) , the \(\varepsilon\) -pseudo-spectrum \(\mathrm{PSp}_{\varepsilon}(u)\) is the set of \(\lambda \in \mathbf{C}\) at distance \(< \varepsilon\) from \(\mathrm{Sp}(u)\) .
\end{enumerate}

Suppose that the spectrum of \(u\) is empty. Then \(u\) is injective, and the endomorphism \(u^{-1} = -\mathrm{R}(u,0)\) is normal (cor. 1, c) and a)). Since \(\mathrm{Sp}(u^{-1}) \subset \{0\}\) (prop. 15, a)), we have \(\mathrm{Sp}(u^{-1}) = \{0\}\) (I, p. 26, cor. 1). Since \(u^{-1}\) is normal, this implies that \(u^{-1} = 0\) (I, p. 110, example 1), which is a contradiction since E is not zero.

To prove \(b\) ), one may assume that \(u\) is the multiplication operator \(m_g\) on \(\mathrm{L}^2 (\mathrm{X},\mu)\) , where \(g\) is a continuous function on a locally compact topological space X equipped with a positive measure \(\mu\) (th. 1 of IV, p. 266). Let \(\lambda \in \mathbf{C} - \mathrm{Sp}(u)\) and let \(\delta >0\) be the distance from \(\lambda\) to the spectrum of \(u\) . To prove that \(\| \mathrm{R}(u,\lambda)\| = \delta^{-1}\) , we reduce to the case \(\lambda = 0\) by replacing \(u\) by \(u - \lambda 1_{\mathrm{E}}\) . The real number \(\delta\) is then the distance from 0 to the spectrum of \(m_g\) . Since the latter coincides with the \(\mu\)-essential image of \(g\) , the result is a consequence of the lemma 3 of IV, p. 182.

Finally, the assertion \(c\) ) follows from \(b\) ) and the definition of \(\mathrm{PSp}_{\varepsilon}(u)\) (IV, p. 250, def. 9).

COROLLARY 3. --- Let u be a partial normal operator on a Hilbert space E. Let \(f \in \mathcal{L}_{\mathrm{u}}(\mathrm{Sp}(u))\) . For all x and y in E, the spectral measure of \((x, y)\) relative to \(f(u)\) is the image measure \(f(\mu)\) , where \(\mu\) is the spectral measure of \((x, y)\) relative to u.

In view of the spectral theorem (th. 1 of IV, p. 266), one may suppose that \(u\) is the multiplication operator \(m_g\) on \(\mathrm{L}^2 (\mathrm{X},\mu)\) , where X is a locally compact topological space, \(\mu\) is a positive measure on X and \(g\in \mathcal{C}(\mathrm{X})\) . Let \(f_{1}\) and \(f_{2}\) be in \(\mathcal{L}^2 (\mathrm{X},\mu)\) with classes \(\widetilde{f}_1\) and \(\widetilde{f}_2\) in \(\mathrm{L}^2 (\mathrm{X},\mu)\) . Since \(f(m_g) = m_{f\circ g}\) (prop. 4, c)), the spectral measure of \((\widetilde{f}_1,\widetilde{f}_2)\) relative to \(f(m_g)\) is the image measure \((f\circ g)(\overline{f}_1f_2\cdot \mu)\) (lemma 4, b) of IV, p. 269). This measure is equal to the image measure \(f(g(\overline{f}_1f_2\cdot \mu))\) (INT, V, p. 72, § 6, n° 4, prop. 4, a)), whence the assertion (lemma 4, b) of IV, p. 269).

COROLLARY 4. --- Let \(g \in \mathcal{L}_{\mathrm{u}}(\mathrm{Sp}(f(u)))\) . We have \(g(f(u)) = (g \circ f)(u)\) .

One has \(g \circ f \in \mathcal{L}_{\mathrm{u}}(\mathrm{Sp}(u))\) (lemma 5 of IV, p. 184). For all \(x\) and \(y\) in E, denote by \(\mu_{x,y}^{\prime}\) the spectral measure of \((x,y)\) relative to \(f(u)\) . By the preceding corollary and INT, V, p. 71, § 6, n°2, th. 1, one has \(g \in \mathcal{L}^{2}(\mathrm{Sp}(f(u)), \mu_{x,y}^{\prime})\) if and only if \(g \circ f \in \mathcal{L}^{2}(\mathrm{Sp}(u), \mu_{x,y})\) , so the domain of \(g(f(u))\) is equal to the domain of \((g \circ f)(u)\) . For all \(x \in \mathrm{E}\) and \(y \in \operatorname{dom}(g(f(u)))\) , one then has the formula

\[
\begin{array}{c} \langle x \mid g (f (u)) y \rangle = \int_ {\mathrm{Sp} (f (u))} g \mu_ {x, y} ^ {\prime} \\ = \int_ {\mathrm{Sp} (f (u))} g f (\mu_ {x, y}) = \langle x \mid (g \circ f) y \rangle \end{array}
\]

(loc. cit.), whence \(g(f(u)) = (g \circ f)(u)\) .

PROPOSITION 5. --- Let u be a normal partial operator on a Hilbert space E.

\begin{enumerate}
\def\labelenumi{\alph{enumi})}
\item
  \(If f \in \mathcal{L}_{\mathrm{u}}^{\infty}(\mathrm{Sp}(u)),\) then \(f(u) \in \mathcal{L}(\mathrm{E})\) ;
\item
  \(If f \in \mathcal{L}_{\mathrm{u}}^{\infty}(\mathrm{Sp}(u))\) and \(g \in \mathcal{L}_{\mathrm{u}}(\mathrm{Sp}(u))\) , then \(f(u) \circ g(u) \subset (fg)(u)\) ;
\item
  The map \(f \mapsto f(u)\) from \(\mathcal{L}_{\mathrm{u}}^{\infty}(\mathrm{Sp}(u))\) into \(\mathcal{L}(\mathrm{E})\) is a continuous unital morphism of star algebras. In particular, one has \(\|f(u)\|\leqslant\|f\|_{\infty}\) for \(f\in\mathcal{L}_{\mathrm{u}}^{\infty}(\mathrm{Sp}(u))\) ;
\item
  If \(u \in \mathcal{L}(\mathrm{E})\) , then for all \(f \in \mathcal{L}_{\mathrm{u}}^{\infty}(\mathrm{Sp}(u))\) , the endomorphism \(f(u)\) belongs to the bicommutant of \(u\) in \(\mathcal{L}(\mathrm{E})\) .
\end{enumerate}

Let \(f \in \mathcal{L}_{\mathrm{u}}^{\infty}(\mathrm{Sp}(u))\) . We have \(\mathrm{D}_f = \mathrm{E}\) and \(f(u)\) is a continuous endomorphism of \(\mathrm{E}\) by formula (8), p. 272, whence the assertion \(a\) ).

Let \(g \in \mathcal{L}_{\mathrm{u}}(\mathrm{Sp}(u))\) Let \(y \in \mathrm{D}_g\) . We have \(y \in \mathrm{D}_{fg}\) and, for all \(x \in \mathrm{E}\) , it follows that \(\langle \overline{f}(u)x | g(u)y \rangle = \langle x | (fg)(u)y \rangle\) by formula (7), p. 271, whence \(f(u)(g(u)y) = (fg)(u)(y)\) , which proves \(b\) .

By \(a\) ) and \(b\) ), the map \(f \mapsto f(u)\) of \(\mathcal{L}_{\mathrm{u}}^{\infty}(\mathrm{Sp}(u))\) in \(\mathcal{L}(\mathrm{E})\) is a unital morphism of involutive algebras; consequently, it is a continuous morphism of norm \(\leqslant 1\) (I, p. 104, prop. 2), whence the assertion \(c\) ).

Suppose that \(u\) is an endomorphism of E. Let \(v \in \mathcal{L}(\mathrm{E})\) commuting with \(u\) . We then have \(v \circ f(u) = f(u) \circ v\) for \(f \in \mathcal{C}(\mathrm{Sp}(u))\) by the properties of the continuous functional calculus (I, p. 110, remark). Let \(x\) and \(y\) be in E. The formulas

\[
\langle x \mid (v \circ f (u)) y \rangle = \langle x \mid (f (u) \circ v) y \rangle ,\tag{9}
\]

valid for every function \(f \in \mathcal{C}(\mathrm{Sp}(u))\) , mean that the spectral measures of \((v^{*}(x), y)\) and of \((x, v(y))\) relative to u are equal. This equality implies, by formula (6), p. 271, that formula (9) is valid for all \(f \in \mathcal{L}_{\mathrm{u}}^{\infty}(\mathrm{Sp}(u))\) . We therefore have \(v \circ f(u) = f(u) \circ v\) .

COROLLARY. --- Let \(f\) and \(g\) in \(\mathcal{L}_{\mathrm{u}}(\mathrm{Sp}(u))\) and \((x,y)\in\mathrm{D}_{f}\times\mathrm{D}_{g}\) . The spectral measure of \((f(u)x,g(u)y)\) relative to \(u\) is the measure \(\overline{f}g\cdot\mu_{x,y}\) , where \(\mu_{x,y}\) is the spectral measure of \((x,y)\) relative to \(u\) .

Denote by \(\nu\) the spectral measure of \((f(u)x,g(u)y)\) relative to \(u\) . For every function \(\varphi \in \mathcal{H}(\mathrm{Sp}(u))\) , we have

\[
\begin{array}{r l} & {\int_ {\mathrm{Sp} (u)} \varphi \nu = \langle f (u) x | \varphi (u) (g (u) y) \rangle} \\ & {\qquad = \langle f (u) x | (\varphi g) (u) y \rangle = \int_ {\mathrm{Sp} (u)} \overline {{f}} \varphi g \mu_ {x, y},} \end{array}
\]

(prop. 5, b)) whence \(\nu = \overline{f} g \cdot \mu_{x,y}\) .

PROPOSITION 6. --- Let u be a normal partial operator on a Hilbert space E. Let \((f_n)_{n \in \mathbf{N}}\) be a sequence in \(\mathcal{L}_{\mathrm{u}}(\mathrm{Sp}(u))\) which converges pointwise to \(f \in \mathcal{L}_{\mathrm{u}}(\mathrm{Sp}(u))\) and such that there exists \(g \in \mathcal{L}_{\mathrm{u}}(\mathrm{Sp}(u))\) satisfying \(|f_n| \leqslant g\) for all \(n \in \mathbf{N}\) . We then have \(\operatorname{dom}(g(u)) \subset \operatorname{dom}(f_n(u))\) for all \(n \in \mathbf{N}\) and \(\operatorname{dom}(g(u)) \subset \operatorname{dom}(f(u))\) . Moreover, for every element y of the domain of \(g(u)\) , we have

\[
f (u) (y) = \lim _ {n \to + \infty} f _ {n} (u) (y).
\]

In particular, if \(f_{n} \in \mathcal{L}_{\mathrm{u}}^{\infty}(\mathrm{Sp}(u))\) for all \(n \in \mathbf{N}\) and if the functions \(f_{n}\) are uniformly bounded, then \(f_{n}(u)\) converges to \(f(u)\) in the space \(\mathcal{L}(\mathrm{E})\) endowed with the topology of pointwise convergence.

Denote by \(\mu_{x,y}\) (resp. \(\mu_{x}\) ) the spectral measure of \((x,y)\) (resp. of x) relative to u.

Let \(y \in \mathrm{dom}(g(u))\) , so that \(g \in \mathcal{L}^2(\mathrm{Sp}(u), \mu_y)\) . The condition \(|f_n| \leqslant g\) implies that \(f_n \in \mathcal{L}^2(\mathrm{Sp}(u), \mu_y)\) , hence \(y \in \mathrm{dom}(f_n(u))\) .

Since \((f_{n})\) converges simply to f and that \(|f_{n}| \leqslant g\) , we have \(f \in \mathcal{L}^{2}(\mathrm{Sp}(u), \mu_{y})\) by the Lebesgue theorem (INT, IV, p. 137, § 3, n° 7, th. 6), hence \(y \in \mathrm{dom}(f(u))\) . Moreover, the sequence \((f_{n})\) converges to f in \(\mathcal{L}^{2}(\mathrm{Sp}(u), \mu_{y})\) , so the norm of \(f_{n}(u)y\) converges to the norm of \(f(u)y\) .

Let \(x \in \mathrm{E}\) . The functions \(f\) and \(g\) , as well as the functions \(f_n\) for all \(n \in \mathbf{N}\) , belong to \(\mathcal{L}^1(\mathrm{Sp}(u), \mu_{x,y})\) (prop. 3). By the Lebesgue theorem (INT, IV, loc. cit.), the sequence \((f_n)\) converges to \(f\) in \(\mathcal{L}^1(\mathrm{Sp}(u), \mu_{x,y})\) , whence

\[
\langle x \mid f _ {n} (u) y \rangle = \int_ {\operatorname{Sp} (u)} f _ {n} \mu_ {x, y} \rightarrow \int_ {\operatorname{Sp} (u)} f \mu_ {x, y} = \langle x \mid f (u) y \rangle .
\]

One concludes that \(f_{n}(u)(y)\) converges to \(f(u)(y)\) (EVT, V, p. 17, prop. 10).

PROPOSITION 7. --- Let X be a locally compact topological space and let ν be a measure on X. Let g: C × X → C be a continuous function with compact support. For z ∈ C, set

\[
h (z) = \int_ {\mathrm{X}} g (z, x) d \nu (x).
\]

\begin{enumerate}
\def\labelenumi{\alph{enumi})}
\item
  The map h from C to C is continuous and bounded;
\item
  The map from X to \(\mathcal{L}(\mathrm{E})\) defined by \(x\mapsto g(u,x)\) is \(\nu\) -integrable and we have
\end{enumerate}

\[
h (u) = \int_ {\mathrm{X}} g (u, x) d \nu (x).
\]

The function h is bounded because g is continuous with compact support, and it is continuous by INT, IV, p. 144, § 4, n° 3, cor. 1. Since g is continuous and has compact support, the map \(x \mapsto g(u, x)\) is continuous from X into \(\mathcal{L}(\mathrm{E})\) . This map has compact support, hence is bounded and integrable on X with respect to \(\nu\) . Denote by

\[
v = \int_ {\mathrm{X}} g (u, x) d \nu (x) \in \mathscr {L} (\mathrm{E}).
\]

Let y and z be elements of E, and let \(\mu\) be the spectral measure of \((y, z)\) relative to u. We have

\[
\langle y \mid v (z) \rangle = \int_ {\mathrm{X}} \langle y \mid g (u, x) z \rangle d \nu (x) = \int_ {\mathrm{X}} \left(\int_ {\operatorname{Sp} (u)} g (\lambda , x) d \mu (\lambda)\right) d \nu (x)
\]

(formula (6), p. 271). Since \(g \in \mathcal{K}(\mathbf{C} \times \mathbf{X})\) , it follows that

\[
\langle y \mid v (z) \rangle = \int_ {\operatorname{Sp} (u)} \Bigl (\int_ {\mathrm{X}} g (\lambda , x) d \nu (x) \Bigr) d \mu (\lambda) = \langle y \mid h (u) z \rangle
\]

by INT, III, p. 84, § 4, n° 1, th. 2 and formula (6), p. 271. This proves that \(v = h(u)\) , as desired.

\subsection{Spectral projectors}

Let \(u\) be a normal partial operator on a complex Hilbert space E. Let A be a universally measurable subset of \(\mathrm{Sp}(u)\) and \(\varphi_{\mathrm{A}}\) its characteristic function. Since \(\varphi_{\mathrm{A}}\) is bounded and satisfies \(\varphi_{\mathrm{A}}^{2} = \varphi_{\mathrm{A}}\) , the endomorphism \(\varphi_{\mathrm{A}}(u)\) of E is an orthogonal projector of E. It is called the spectral projector of \(u\) defined by A. We denote \(p_{\mathrm{A}} = \varphi_{\mathrm{A}}(u)\) . We have \(p_{\emptyset} = 0\) and \(p_{\mathrm{Sp}(u)} = 1_{\mathrm{E}}\) .

Let A be a universally measurable subset of C. The spectral projector of u defined by A is the spectral projector \(p_{\mathrm{Sp}(u)\cap\mathrm{A}}\) . It is also denoted simply \(p_{A}\) . For every function \(f \in \mathcal{L}_{\mathrm{u}}(\mathrm{Sp}(u))\) , for all \(x \in E\) and every \(y \in \mathrm{dom}(f(u))\) , we have the formula

\[
\langle p _ {\mathrm{A}} (x) \mid f (u) y \rangle = \int_ {\operatorname{Sp} (u) \cap \mathrm{A}} f d \mu ,\tag{10}
\]

where \(\mu\) is the spectral measure of \((x,y)\) relative to \(u\) (formula (7), p. 271). Let A and B be universally measurable subsets of \(\mathrm{Sp}(u)\) . Since \(\varphi_{\mathrm{A}}\varphi_{\mathrm{B}} = \varphi_{\mathrm{A}\cap \mathrm{B}}\) , we have \(p_{\mathrm{A}} \circ p_{\mathrm{B}} = p_{\mathrm{A}\cap \mathrm{B}}\) (prop. 5 of IV, p. 275, \(c\) )). In particular, if A and B are disjoint, the images of \(p_{\mathrm{A}}\) and of \(p_{\mathrm{B}}\) are orthogonal.

PROPOSITION 8. --- Let \((\mathrm{A}_{i})_{i\in\mathrm{I}}\) be a countable family of pairwise disjoint universally measurable subsets of \(\mathrm{Sp}(u)\) and let \(p_{i}\) be the spectral projector of u defined by \(A_{i}\) . The union A of the sets \(A_{i}\) is a universally measurable subset of \(\mathrm{Sp}(u)\) and the series \(\sum_{i}p_{i}\) converges to \(p_{A}\) in \(\mathcal{L}(\mathrm{E})\) endowed with the topology of pointwise convergence.

The set A is universally measurable by INT, IV, p. 177, § 5, n° 4, cor. 2. The series \(\sum_{i} \varphi_{\mathrm{A}_{i}}\) converges simply to \(\varphi_{\mathrm{A}}\) and its partial sums are bounded by 1. The assertion therefore follows from prop. 6 of IV, p. 276.

PROPOSITION 9. --- Let A be a closed subset of Sp(u). Let \(\varphi_{\mathrm{A}}\) be the characteristic function of A and \(p_{\mathrm{A}} = \varphi_{\mathrm{A}}(u)\) the spectral projector of u defined by A. Denote by \(\mathrm{E}_{\mathrm{A}}\) the image of \(p_{\mathrm{A}}\) . This is a closed subspace of E.

\begin{enumerate}
\def\labelenumi{\alph{enumi})}
\item
  The subspace \(E_{A}\) is the set of \(x \in E\) such that the support of the spectral measure of x relative to u is contained in A;
\item
  If A is bounded in \(\mathbf C\), then \(E_{A}\) is contained in the domain of u;
\item
  For all x belonging to the domain of u, we have \(p_{\mathrm{A}}(x) \in \operatorname{dom}(u)\) and \(u(p_{\mathrm{A}}(x)) \in \mathrm{E}_{\mathrm{A}}\) in particular \(u(x) \in \mathrm{E}_{\mathrm{A}}\) if \(x \in \operatorname{dom}(u) \cap \mathrm{E}_{\mathrm{A}}\) .
\end{enumerate}

For \(x\) in E, we denote by \(\mu_x\) the spectral measure of \(x\) relative to \(u\) .

Let us prove \(a\) ). Let \(x\) be an element of \(\mathrm{E}_{\mathrm{A}}\) . Let \(z \in \mathbf{C} - \mathbf{A}\) and let U be a relatively compact open neighborhood of \(z\) which does not meet A. For every function \(f \in \mathcal{K}(\mathbf{C})\) with support contained in U, we have

\[
\int_ {\mathrm{Sp} (u)} f \mu_ {x} = \langle x | f (u) x \rangle = \langle p _ {\mathrm{A}} (x) | f (u) x \rangle = \int_ {\mathrm{Sp} (u) \cap \mathrm{A}} f \mu_ {x} = 0
\]

by formula (10). This means that \(z\) does not belong to the support of \(\mu_x\) . Consequently, the support of \(\mu_x\) is contained in A.

Conversely, let \(x \in E\) such that the support of \(\mu_{x}\) is contained in A. It follows that

\[
\langle x \mid p _ {\mathrm{A}} (x) \rangle = \int_ {\operatorname{Sp} (u) \cap \mathrm{A}} \mu_ {x} = \int_ {\operatorname{Sp} (u)} \mu_ {x} = \langle x \mid x \rangle ,
\]

\((loc. \ cit.) \ \text{therefore} \ \|p_{\mathrm{A}}(x)-x\|^{2} = \|p_{\mathrm{A}}(x)\|^{2} - \|x\|^{2} \leqslant 0 \ \text{and consequently } p_{\mathrm{A}}(x) = x, \ \text{that is, } x \in \mathrm{E}_{\mathrm{A}}.\)

Assertion b) follows from a) and the remark in IV, p. 273.

Let us prove \(c\) ). The domain of \(u\) is the set of \(x \in \mathrm{E}\) such that the identity function of \(\operatorname{Sp}(u)\) belongs to \(\mathcal{L}^2(\operatorname{Sp}(u), \mu_x)\) (loc. cit.). Since \(\mu_{p_{\mathrm{A}}(x)} = \varphi_{\mathrm{A}} \cdot \mu_x\) for \(x \in \mathrm{E}\) (corollary of Prop. 5 of IV, p. 275), we have \(p_{\mathrm{A}}(x) \in \operatorname{dom}(u)\) if \(x \in \operatorname{dom}(u)\) .

Let \(x \in \operatorname{dom}(u)\) and \(y = p_{\operatorname{A}}(x)\) ; we have \(y \in \operatorname{dom}(u) \cap \operatorname{E}_{\operatorname{A}}\) according to \(c\) ). The spectral measure of \(u(y)\) relative to \(u\) is \(|\operatorname{Id}_{\operatorname{Sp}(u)}|^2 \cdot \mu_y\) (loc. cit.). Since \(\mu_y\) is supported in A, the same holds for \(\mu_{u(y)}\) , hence \(u(y) \in \operatorname{E}_{\operatorname{A}}\) according to \(a\) .

COROLLARY. --- Let \(\lambda \in \mathrm{Sp}(u)\) The image of \(p_{\{\lambda\}}\) is the eigensubspace of \(u\) relative to \(\lambda\) .

Let \(x \in \operatorname{dom}(u)\) . Denote by \(\mu\) (resp. \(\nu\) ) the spectral measure of the element \(x\) relative to \(u\) (resp. the spectral measure of \((x, u(x))\) relative to \(u\) ). We have \(\nu = \operatorname{Id}_{\operatorname{Sp}(u)} \cdot \mu\) (corollary of Prop. 5 of IV, p. 275). If \(u(x) = \lambda x\) , we also have \(\nu = \lambda \mu\) , whence the equality \(\operatorname{Id}_{\operatorname{Sp}(u)} \cdot \mu = \lambda \mu\) . This implies that the support of \(\mu\) is contained in \(\{\lambda\}\) (cf. INT, V, p. 46, § 5, no. 3, cor. 2) and therefore \(x\) belongs to the image of \(p_{\{\lambda\}}\) (prop. 9, a)).

Conversely, suppose that \(x\) belongs to the image \(\mathrm{E}_{\lambda}\) of \(p_{\{\lambda\}}\) . Then \(x\) belongs to \(\operatorname{dom}(u)\) and \(u(x)\) also belongs to \(\mathrm{E}_{\lambda}\) (loc. cit., \(b\) )). Since \(\varphi_{\{\lambda\}} \cdot (\operatorname{Id}_{\operatorname{Sp}(u)} - \lambda) = 0\) , we have the relation \(p_{\{\lambda\}} \circ (u - \lambda 1_{\mathrm{E}}) \subset 0\) (prop. 5 of IV, p. 275, \(c\) )), whence \(0 = p_{\{\lambda\}}(u(x)) - \lambda p_{\{\lambda\}}(x) = u(x) - \lambda x\) .

PROPOSITION 10. — Let \(\lambda \in \mathbf{C}\) . We have \(\lambda \in \mathrm{Sp}(u)\) if and only if, for every open neighborhood V of \(\lambda\) in \(\mathbf{C}\) , the spectral projector \(p_{\mathrm{V}}\) of u relative to V is nonzero.

If \(\lambda \notin \mathrm{Sp}(u)\) , then there exists an open neighborhood \(V\) of \(\lambda\) in \(\mathbf{C}\) which does not meet \(\mathrm{Sp}(u)\) , and then \(p_{\mathrm{V}} = p_{\varnothing} = 0\) .

Conversely, suppose that there exists an open neighborhood V of \(\lambda\) in \(\mathbf{C}\) such that \(p_{V} = 0\) . Let c > 0 be such that the disk with center \(\lambda\) and radius c is contained in V.

Let \(x \in \mathrm{dom}(u)\) and let \(\mu_x\) the spectral measure of \(x\) relative to \(u\) . Since \(\mu_x(\mathrm{V}) = \langle x | p_{\mathrm{V}}(x) \rangle = 0\) , one computes

\[
\begin{array}{r l} & {\int_ {\mathbf {C}} | z - \lambda | ^ {2} d \mu_ {x} (z) = \int_ {\mathbf {C - V}} | z - \lambda | ^ {2} d \mu_ {x} (z)} \\ & {\qquad \geqslant c ^ {2} \int_ {\mathbf {C - V}} d \mu_ {x} (z) = c ^ {2} \int_ {\mathbf {C}} d \mu_ {x} (z) = c ^ {2} \| x \| ^ {2}.} \end{array}
\]

But, on the other hand, we have

\[
\| u (x) - \lambda x \| ^ {2} = \int_ {\mathbf {C}} | z - \lambda | ^ {2} d \mu_ {x} (z) = \| u ^ {*} (x) - \overline {{\lambda}} x \| ^ {2}
\]

(formula (8), p. 272), whence \(\|u(x)-\lambda x\|\geqslant c\|x\|\) and \(\|u^{*}(x)-\overline{\lambda}x\|\geqslant c\|x\|\) . It follows that \(\lambda\) belongs to the resolvent set of u (lemma 5 of IV, p. 248). This concludes the proof.

COROLLARY. — Let A be an open set in \(\mathbf{C}\) and \(n \in \mathbf N\) . If A contains n elements of \(\mathrm{Sp}(u)\) , then the dimension of the image of the spectral projector \(p_{A}\) of u relative to A is at least equal to n.

Let \(\lambda_1, \ldots, \lambda_n\) be distinct elements of the spectrum of \(u\) belonging to A. There exists a family \((\mathrm{V}_i)_{1 \leqslant i \leqslant n}\) of pairwise disjoint open sets of \(\mathbf{C}\) such that \(\lambda_i \in \mathrm{V}_i\) for \(1 \leqslant i \leqslant n\) . Let B be the union of the sets \(\mathrm{V}_i\) . The image of \(p_{\mathrm{A}}\) contains the image of the spectral projector \(p_{\mathrm{B}}\) ; moreover, since \(p_{\mathrm{B}}\) is the sum of the projectors \(p_{\mathrm{V}_i}\) , and since the image of \(p_{\mathrm{V}_i}\) is orthogonal to the image of \(p_{\mathrm{V}_j}\) for all \(i \neq j\) , the result follows from Prop. 10.

\subsection{The Helffer--Sjöstrand formula}

In this issue, E denotes a complex Hilbert space. We shall obtain a formula for certain cases of the functional calculus of a partial self-adjoint operator, expressed directly in terms of the resolvent of the operator in question.

We endow R (resp. C) with Lebesgue measure, denoted \(\mu\) and we identify the group R with its dual group by the map \((x,y)\mapsto\exp(2i\pi xy)\) (cf. Corollary 3 of II, p. 236).

For any function \(f\) defined and differentiable on an open set U of \(\mathbf{R}^{2}\) , identified with \(\mathbf{C}\) with real coordinates \(x\) and \(y\) , we denote

\[
{\frac {\partial f}{\partial \overline {{{z}}}}} = {\frac {1}{2}} {\Bigl (} {\frac {\partial f}{\partial x}} + i {\frac {\partial f}{\partial y}} {\Bigr)}
\]

(cf. VAR, R2, 8.8.10, p. 24).

Lemma 5. --- Let \(f \in \mathcal{D}(\mathbf{R})\) . There exists a function \(\tilde{f}\) in \(\mathcal{D}(\mathbf{C})\) which coincides with \(f\) on \(\mathbf{R}\) and satisfies

\[
\frac {\partial \widetilde {f}}{\partial \overline {{z}}} (x, 0) = 0\tag{11}
\]

for all \(x \in \mathbf{R}\) . We then have

\[
\frac {\partial \widetilde {f}}{\partial y} (x, 0) = i f ^ {\prime} (x)\tag{12}
\]

for all \(x \in R\) and there exists a real number C ≥ 0 such that

\[
\left| \frac {\partial \widetilde {f}}{\partial \overline {{z}}} (x, y) \right| \leqslant C | y |\tag{13}
\]

for all \((x,y)\in \mathbf{R}^2\)

There exists \(\varphi\in\mathcal{D}(\mathbf{R})\) whose support is contained in \([-2,2]\) and which is equal to 1 on \([-1,1]\) (lemma 1 of IV, p. 196). Set

\[
\widetilde {f} (x, y) = \bigl (f (x) + i y f ^ {\prime} (x) \bigr) \varphi (y)
\]

for \((x,y)\in\mathbf{R}^{2}\) . We have \(\widetilde{f}\in\mathcal{D}(\mathbf{C})\) and \(\widetilde{f}\) coincides with f on R. Moreover, for any \((x,y)\in\mathbf{R}^{2}\) , it follows that

\[
\frac {\partial \widetilde {f}}{\partial \overline {{z}}} (x, y) = \frac {1}{2} \Big ((i f (x) - y f ^ {\prime} (x)) \varphi^ {\prime} (y) + i y f ^ {\prime \prime} (x) \varphi (y) \Big).
\]

As the function \(\varphi\) is equal to 1 in a neighborhood of 0, we have \(\varphi'(0) = 0\) whence (11).

Let \(\tilde{f}\) be in \(\mathcal{D}(\mathbf{C})\) satisfying (11). Formula (12) follows from this, and the bound (13) is obtained using the mean value theorem (FVR, I, p. 23, th. 2).

We say that \(\tilde{f}\) is an almost analytic extension of f.

Lemma 6. --- Let \(\varepsilon > 0\) . The function \(\sigma_{\varepsilon}\) defined on \(\mathbf{R}\) by

\[
\sigma_ {\varepsilon} (x) = \frac {2 i \varepsilon x}{x ^ {2} + \varepsilon^ {2}}
\]

belongs to \(L^{2}(\mathbf{R})\) Its Fourier transform is the class in \(L^{2}(\mathbf{R})\) of the function \(\eta_{\varepsilon}\) which vanishes at 0 and satisfies

\[
\eta_ {\varepsilon} (y) = \frac {2 \pi \varepsilon y}{| y |} e ^ {- 2 \pi \varepsilon | y |}
\]

for all \(y \neq 0\) .

One has \(\sigma_{\varepsilon} \in \mathrm{L}^{2}(\mathbf{R})\) by prop. 3 of IV, p. 199, and the class of the function \(\eta_{\varepsilon}\) belongs to \(\mathrm{L}^{2}(\mathbf{R}) \cap \mathrm{L}^{1}(\mathbf{R})\) . For every \(x \in \mathbf{R}\) , we have

\[
\begin{array}{r} \overline {{\mathcal {F}}} (\eta_ {\varepsilon}) (x) = 2 \pi \varepsilon \int_ {\mathbf {R} _ {+}} e ^ {2 \pi (i x - \varepsilon) y} d y - 2 \pi \varepsilon \int_ {\mathbf {R} _ {-}} e ^ {2 \pi (i x + \varepsilon) y} d y \\ = \frac {\varepsilon}{\varepsilon - i x} - \frac {\varepsilon}{\varepsilon + i x} = \frac {2 i \varepsilon x}{x ^ {2} + \varepsilon^ {2}}, \end{array}
\]

whence \(\overline{\mathcal{F}}(\eta_{\varepsilon}) = \sigma_{\varepsilon}\) The result then follows from the Fourier inversion formula in \(\mathrm{L}^2(\mathbf{R})\) (corollary of Theorem 2 of II, p. 220).

Lemma 7. --- Let \(f \in \mathcal{D}(\mathbf{R})\) and let \(\widetilde{f} \in \mathcal{D}(\mathbf{C})\) be an almost analytic extension of \(f\) . For \(\varepsilon > 0\) , define \(f_{\varepsilon}\) on \(\mathbf{R}\) by

\[
f _ {\varepsilon} (x) = - \frac {1}{2 i \pi} \int_ {\mathbf {R}} \biggl (\frac {\widetilde {f} (y + i \varepsilon)}{y - x + i \varepsilon} - \frac {\widetilde {f} (y - i \varepsilon)}{y - x - i \varepsilon} \biggr) d y.
\]

Then \(f_{\varepsilon}\) is continuous and bounded, and \(f_{\varepsilon}\) converges to \(f\) in \(\mathcal{C}_b(\mathbf{R})\) as \(\varepsilon\) tends to 0.

The continuity of \(f_{\varepsilon}\) is a consequence of INT, IV, p. 144, § 4, n° 3, cor. 1, since \(\widetilde{f}\) is compactly supported. Moreover, if \(r\) is such that the support of \(\widetilde{f}\) is contained in \([-r, r] \times \mathbf{R}\) , we have

\[
| f _ {\varepsilon} (x) | \leqslant 2 \| \widetilde {f} \| _ {\infty} \int_ {- r} ^ {r} \frac {1}{\sqrt {(x - y) ^ {2} + \varepsilon^ {2}}} d y \leqslant \frac {4 r}{\varepsilon} \| \widetilde {f} \| _ {\infty},
\]

hence \(f_{\varepsilon}\) is bounded.

The first-order Taylor expansion (FVR, I, p. 30, prop. 3) and formula (13) show that there exist M ≥ 0 and a function \(\varrho_{1}\) with values in \(R^{2}\) such that

\[
\widetilde {f} (y + i \gamma) = f (y) + i \gamma f ^ {\prime} (y) + \gamma^ {2} \varrho_ {1} (y; \gamma), \text {   and   } | \varrho_ {1} (y; \gamma) | \leqslant M,
\]

for all \(y \in \mathbf{R}\) and \(\gamma \in \mathbf{R}\) . Since the map \(y \mapsto \widetilde{f}(y + i\gamma)\) has support contained in \([-r, r]\) for all \(\gamma \in \mathbf{R}\) , the map \(y \mapsto \varrho_{1}(y; \gamma)\) has support contained in \([-r, r]\) , for all \(\gamma \in \mathbf{R}\) .

Let \(g_{\varepsilon}\) the function defined on \(\mathbf{R}^2\) by

\[
g _ {\varepsilon} (x, y) = \frac {\widetilde {f} (y + i \varepsilon)}{y - x + i \varepsilon} - \frac {\widetilde {f} (y - i \varepsilon)}{y - x - i \varepsilon}.
\]

For every \((x,y)\in \mathbf{R}^2\) and \(\varepsilon >0\) , we obtain

\[
\begin{array}{r l} g _ {\varepsilon} (x, y) = - \frac {2 i \varepsilon}{(x - y) ^ {2} + \varepsilon^ {2}} f (y) + \frac {2 i \varepsilon (y - x)}{(x - y) ^ {2} + \varepsilon^ {2}} f ^ {\prime} (y) \\ & \quad + \varepsilon^ {2} \bigg (\frac {\varrho_ {1} (y ; \varepsilon)}{y - x + i \varepsilon} - \frac {\varrho_ {1} (y ; - \varepsilon)}{y - x - i \varepsilon} \bigg). \end{array}
\]

Let \(x \in \mathbf{R}\) and \(\varepsilon > 0\) . We have

\[
\begin{array}{c} \left| \varepsilon^ {2} \int_ {\mathbf {R}} \left(\frac {\varrho_ {1} (y ; \varepsilon)}{y - x + i \varepsilon} - \frac {\varrho_ {1} (y ; - \varepsilon)}{y - x - i \varepsilon}\right) d y \right| \leqslant 2 \mathrm{M} \varepsilon^ {2} \int_ {- r} ^ {r} \frac {d y}{\sqrt {(x - y) ^ {2} + \varepsilon^ {2}}} \\ \leqslant 4 \mathrm{Mr} \varepsilon . \end{array}
\]

Consequently, for every \(x \in \mathbf{R}\) , it follows that

\[
f _ {\varepsilon} (x) = - \frac {1}{2 i \pi} \int_ {\mathbf {R}} g _ {\varepsilon} (x, y) d y = (f * \delta_ {\varepsilon}) (x) - \frac {1}{2 i \pi} (f ^ {\prime} * \sigma_ {\varepsilon}) (x) + k _ {\varepsilon} (x)
\]

where \(\delta_{\varepsilon}\) and \(\sigma_{\varepsilon}\) are the functions on \(\mathbf{R}\) defined by

\[
\delta_ {\varepsilon} (x) = \frac {1}{\pi} \frac {\varepsilon}{x ^ {2} + \varepsilon^ {2}}, \qquad \sigma_ {\varepsilon} (x) = \frac {2 i \varepsilon x}{x ^ {2} + \varepsilon^ {2}},
\]

and \(\| k_{\varepsilon}\|_{\infty}\leqslant 2\mathrm{Mr}\varepsilon .\)

The function \(\mathcal{F}(f') \in \mathcal{S}(\mathbf{R})\) is integrable (prop. 18 of IV, p. 219 and prop. 13 of IV, p. 213). By lemma 6, it follows that

\[
\begin{array}{r l r} & & {\int_ {\mathbf {R}} | \mathcal {F} (f ^ {\prime}) (y) \mathcal {F} (\sigma_ {\varepsilon}) (y) | d y = 2 \pi \varepsilon \int_ {\mathbf {R}} | \mathcal {F} (f ^ {\prime}) (y) | e ^ {- 2 \pi \varepsilon | y |} d y} \\ & & {\leqslant 2 \pi \varepsilon \int_ {\mathbf {R}} | \mathcal {F} (f ^ {\prime}) (y) | d y,} \end{array}
\]

hence \(\mathcal{F}(f')\mathcal{F}(\sigma_{\varepsilon})\) converges to 0 in \(\mathrm{L}^{1}(\mathbf{R})\) as \(\varepsilon\) tends to 0. Since \(f'\) and \(\sigma_{\varepsilon}\) belong to \(\mathrm{L}^{2}(\mathbf{R})\) , prop. 14 of II, p. 223 implies that \(f'*\sigma_{\varepsilon}=\overline{\mathcal{F}}(\mathcal{F}(f')\mathcal{F}(\sigma_{\varepsilon}))\) converges to 0 in \(\mathcal{C}_{b}(\mathbf{R})\) .

For every \(\varepsilon > 0\) , the positive measure \(\delta_{\varepsilon} \cdot dx\) on \(\mathbf{R}\) has total mass 1 (cf. FVR, III, p. 7). The set of measures \(\delta_{\varepsilon} \cdot dx\) for \(\varepsilon > 0\) and the filter induced on this set by the filter of neighborhoods of 0 in \(\mathbf{R}_{+}^{*}\) satisfy the hypotheses of lemma 4 of INT, VIII, p. 137, § 2, no. 7. The measures \(\delta_{\varepsilon} \cdot dx\) therefore converge in \(\mathcal{M}^{1}(\mathbf{R})\) toward the point measure \(\varepsilon_{0}\) as \(\varepsilon\) tends to 0. It follows that \(f * \delta_{\varepsilon} \to f\) in \(\mathcal{C}_{b}(\mathbf{R})\) (INT, VIII, p. 163, § 4, no. 4). The lemma is proved.

We denote by \(\mu\) the Lebesgue measure on \(\mathbf{C}\).

THEOREM 2 (Helffer--Sjöstrand). --- Let u be a partial self-adjoint operator on E. Let \(f \in \mathcal{D}(\mathbf{R})\) and \(\widetilde{f} \in \mathcal{D}(\mathbf{C})\) be an almost analytic extension of \(f\) . Let h be the map from \(\mathbf{C}\) into \(\mathcal{L}(\mathrm{E})\) defined by

\[
h (\lambda) = \frac {\partial \widetilde {f}}{\partial \overline {{z}}} (\lambda) \mathrm{R} (u, \lambda)
\]

if \(\lambda\in\mathbf{C} - \mathbf{R}\) and \(h(\lambda)=0\) if \(\lambda\in\mathbf{R}\) . Then h is \(\mu\) -integrable on \(\mathbf{C}\) and

\[
f (u) = - \frac {1}{\pi} \int_ {\mathbf {C}} h (\lambda) d \mu (\lambda).
\]

The map h is measurable and its support is compact. Since u is self-adjoint, we have \(\|\mathrm{R}(u,\lambda)\| \leqslant |\mathcal{I}(\lambda)|^{-1}\) for all \(\lambda \in \mathbf{C} - \mathbf{R}\) (prop. 17 of IV, p. 248). The map h is therefore bounded according to formula (13), and consequently it is integrable on \(\mathbf{C}\).

Let \(\varepsilon \in \mathbf{R}_{+}^{*}\) . We denote \(\mathrm{F}_{\varepsilon}^{+}\) (resp. \(\mathrm{F}_{\varepsilon}^{-}\) ) the closed set in \(\mathbf{C}\) of the \(\lambda \in \mathbf{C}\) such that \(\mathcal{I}(\lambda) \geqslant \varepsilon\) (resp. \(\mathcal{I}(\lambda) \leqslant -\varepsilon\) ). We have

\[
\int_ {\mathbf {C}} h (\lambda) d \mu (\lambda) = \lim _ {\varepsilon \to 0} \Bigl (\int_ {\mathrm{F} _ {\varepsilon} ^ {+}} h (\lambda) d \mu (\lambda) + \int_ {\mathrm{F} _ {\varepsilon} ^ {-}} h (\lambda) d \mu (\lambda) \Bigr).
\]

Let r > 0 such that the support of h is contained in \(\mathbf{C} = [-r, r]^{2}\) . Denote by \(R_{\varepsilon}^{+}\) the box \([-2r, 2r] \times [\varepsilon, 2r]\) in \(\mathbf{C}\). It is a locally polyhedral subset of \(\mathbf{C}\) (VAR, R2, 11.3, p. 48). It satisfies the following conditions:

\begin{enumerate}
\def\labelenumi{(\roman{enumi})}
\item
  One has \(R_{\varepsilon}^{+} \subset 2C \cap F_{\varepsilon}^{+}\) ;
\item
  The set \(R_{\varepsilon}^{+}\) contains the intersection of \(F_{\varepsilon}^{+}\) and the support of h;
\item
  The regular boundary \(\partial R_{\varepsilon}^{+}\) (VAR, R2, 11.3.2, p. 49) contains the segment \(S_{\varepsilon} = [-r, r] + i\varepsilon \subset C\) ;
\item
  One has \(h(\lambda) = 0\) if \(\lambda \in \partial \mathrm{R}_{\varepsilon}^{+} - \mathrm{S}_{\varepsilon}\) .
\end{enumerate}

Denote by \(d\lambda\) (resp. \(d\overline{\lambda}\) ) the differential 1-form on \(\mathbf{C}\) of the identity map of \(\mathbf{C}\) (resp. of complex conjugation). Denote by \(g\) the function on \(\mathbf{C} - \mathbf{R}\) with values in \(\mathcal{L}(\mathrm{E})\) such that \(g(\lambda) = \widetilde{f}(\lambda)\mathrm{R}(u,\lambda)\) for \(\lambda \in \mathbf{C} - \mathbf{R}\) . Let \(\omega = g d\lambda\) ; it is a differential form of degree 1 on \(\mathbf{C} - \mathbf{R}\) with compact support and values in \(\mathcal{L}(\mathrm{E})\) . Since the resolvent of \(u\) is holomorphic (prop. 14 of IV, p. 246), we have

\[
d \omega = \left(\frac {\partial \widetilde {f}}{\partial \overline {{z}}} (\lambda) \mathrm{R} (u, \lambda) + \widetilde {f} (\lambda) \frac {\partial}{\partial \overline {{z}}} \mathrm{R} (u, \lambda)\right) d \overline {{\lambda}} \wedge d \lambda = - h (\lambda) d \lambda \wedge d \overline {{\lambda}}.
\]

The vector measure associated with \(d\omega\) (VAR, R2, 10.4.3, p. 43) is the density measure -2ih with respect to Lebesgue measure. Applying Stokes' formula to the locally polyhedral set \(R_{\varepsilon}^{+}\) and to the differential form \(\omega\) (VAR, R2, 11.3.4, p. 49), we thus obtain

\[
\frac {i}{2} \int_ {\mathrm{R} _ {\varepsilon} ^ {+}} d \omega = \frac {i}{2} \int_ {\partial \mathrm{R} _ {\varepsilon} ^ {+}} \omega = \frac {i}{2} \int_ {\mathrm{S} _ {\varepsilon}} \omega = \frac {i}{2} \int_ {- r} ^ {r} \widetilde {f} (y + i \varepsilon) \mathrm{R} (u, y + i \varepsilon) d y,
\]

whence

\[
\int_ {\mathrm{F} _ {\varepsilon} ^ {+}} h (\lambda) d \mu (\lambda) = - \frac {i}{2} \int_ {\mathbf {R}} \widetilde {f} (y + i \varepsilon) \mathrm{R} (u, y + i \varepsilon) d y.
\]

Reasoning in the same way for \(F_{\varepsilon}^{-}\) , we obtain

\[
\int_ {\mathrm{F} _ {\varepsilon} ^ {-}} h (\lambda) d \mu (\lambda) = - \frac {i}{2} \int_ {\mathbf {R}} \widetilde {f} (y - i \varepsilon) \mathrm{R} (u, y - i \varepsilon) d y,
\]

and we conclude that the integral of \(h\) over \(\mathbf{C}\) is the limit when \(\varepsilon\to0\) of

\[
v _ {\varepsilon} = \frac {i}{2} \int_ {\mathbf {R}} \Bigl (\widetilde {f} (y + i \varepsilon) \mathrm{R} (u, y + i \varepsilon) - \widetilde {f} (y - i \varepsilon) \mathrm{R} (u, y - i \varepsilon) \Bigr) d y.
\]

By prop. 7, we have \(v_{\varepsilon} = \pi f_{\varepsilon}(u)\) , where \(f_{\varepsilon}\) is the function defined on \(\mathbf{R}\) by

\[
f _ {\varepsilon} (x) = - \frac {1}{2 i \pi} \int_ {\mathbf {R}} \biggl (\frac {\widetilde {f} (y + i \varepsilon)}{y - x + i \varepsilon} - \frac {\widetilde {f} (y - i \varepsilon)}{y - x - i \varepsilon} \biggr) d y.
\]

Since \(f_{\varepsilon} \to f\) as \(\varepsilon \to 0\) in \(\mathcal{C}_{b}(\mathbf{R})\) (Lemma 7), the endomorphism \(v_{\varepsilon} = \pi f_{\varepsilon}(u)\) converges to \(\pi f(u)\) in \(\mathcal{L}(\mathrm{E})\) (prop. 5 of IV, p. 275). The theorem is proved.

\subsection{Resolvent topologies and continuity of the functional calculus}

In this issue, E denotes a complex Hilbert space. Recall that we denote by \(\mathcal{A}(\mathrm{E})\) the set of partial self-adjoint operators on E. We will extend to \(\mathcal{A}(\mathrm{E})\) the continuity properties of item 8 of IV, p. 194.

DEFINITION 6. --- Let T be a locally convex topology on \(\mathcal{L}(\mathrm{E})\) . The T-resolvent topology on \(\mathcal{A}(\mathrm{E})\) is the coarsest topology such that the maps \(u \mapsto \mathrm{R}(u, \lambda)\) from \(\mathcal{A}(\mathrm{E})\) into \(\mathcal{L}(\mathrm{E})\) endowed with the topology T are continuous for every \(\lambda \in \mathbf{C} - \mathbf{R}\) .

PROPOSITION 11. --- Let T be a locally convex topology on \(\mathcal{L}(\mathrm{E})\) which is coarser than the Banach space topology of \(\mathcal{L}(\mathrm{E})\) .

\begin{enumerate}
\def\labelenumi{\alph{enumi})}
\item
  Let \(f \in \mathcal{C}_0(\mathbf{R})\) . For every sequence \((u_n)_{n \in \mathbf{N}}\) in \(\mathcal{A}(\mathrm{E})\) which converges to u for the topology \(\mathcal{T}\) -resolvent, the sequence \((f(u_n))_{n \in \mathbf{N}}\) converges to \(f(u)\) in the space \(\mathcal{L}(\mathrm{E})\) endowed with the topology \(\mathcal{T}\) ;
\item
  Suppose that everything \(u \in \mathcal{A}(\mathrm{E})\) admits a countable fundamental system of neighborhoods for the topology \(\mathcal{T}\) -resolvent. The map of \(\mathcal{C}_0(\mathbf{R}) \times \mathcal{A}(\mathrm{E})\) in the space \(\mathcal{L}(\mathrm{E})\) endowed with the topology \(\mathcal{T}\) defined by \((f, u) \mapsto f(u)\) is continuous.
\end{enumerate}

Let us prove \(a\) ). The topology \(\mathcal{T}\) on \(\mathcal{L}(\mathrm{E})\) is the upper bound of the topologies defined by the semi-norms continuous for the topology \(\mathcal{T}\) (EVT, II, p. 26, cor. and following remark). It therefore suffices to prove that for every seminorm \(p\) on \(\mathcal{L}(\mathrm{E})\) which is continuous for the topology \(\mathcal{T}\) the sequence \(p(f(u_n) - f(u))\) converges to 0 (TG, I, p. 51, prop. 10).

Let p be such a seminorm. Denote by \(\mathcal{L}(\mathrm{E})_{p}\) the separated Banach space completion of the space \(\mathcal{L}(\mathrm{E})\) equipped with the seminorm p, and denote by \(\varpi\) the canonical map from \(\mathcal{L}(\mathrm{E})\) into \(\mathcal{L}(\mathrm{E})_{p}\) ; it is continuous, and one has \(p(u) = \|\varpi(u)\|_{p}\) for all \(u \in \mathcal{L}(\mathrm{E})\) , where the norm is that of the space \(\mathcal{L}(\mathrm{E})_{p}\) .

Since T is coarser than the Banach space topology of \(\mathcal{L}(\mathrm{E})\) , there exists \(c \geqslant 0\) such that \(p(u) \leqslant c \|u\|\) for all \(u \in \mathcal{L}(\mathrm{E})\) .

Suppose first that \(f \in \mathcal{D}(\mathbf{R})\) Let \(\tilde{f}\) an almost analytic extension of \(f\) . Denote by \(\mu\) the Lebesgue measure on \(\mathbf{C}\) . By the Helffer--Sjöstrand formula (th. 2 of IV, p. 284), one has

\[
\varpi (f (u _ {n})) = - \frac {1}{\pi} \int_ {\mathbf {C}} \frac {\partial \widetilde {f}}{\partial \bar {z}} (\lambda) \varpi (\mathrm{R} (u _ {n}, \lambda)) d \mu (\lambda)\tag{14}
\]

for all \(n \in \mathbf{N}\) .

Let \(\lambda \in \mathbf{C} - \mathbf{R}\) By hypothesis, the sequence of resolvents \(\mathrm{R}(u_n,\lambda)\) converges to \(\mathrm{R}(u,\lambda)\) in \(\mathcal{L}(\mathrm{E})\) endowed with the topology \(\mathcal{T}\) , hence the sequence \((\varpi (\mathrm{R}(u_n,\lambda)))_n\) converges to \(\varpi (\mathrm{R}(u,\lambda))\) to the Banach space \(\mathcal{L}(\mathrm{E})_p\) .

For every \(n \in \mathbf{N}\) , we have

\[
\left\| \frac {\partial \widetilde {f}}{\partial \bar {z}} (\lambda) \mathrm{R} (u _ {n}, \lambda) \right\| \leqslant \left| \frac {1}{\mathcal{I} (\lambda)} \frac {\partial \widetilde {f}}{\partial \bar {z}} (\lambda) \right|
\]

(prop. 17 of IV, p. 248) therefore

\[
\left\| \frac {\partial \widetilde {f}}{\partial \overline {{z}}} (\lambda) \varpi (\mathrm{R} (u _ {n}, \lambda)) \right\| _ {p} \leqslant c \left| \frac {1}{\mathcal {I} (\lambda)} \frac {\partial \widetilde {f}}{\partial \overline {{z}}} (\lambda) \right|.
\]

The estimate (13) of IV, p. 281 shows that the right-hand side of this inequality is a bounded function for \(\lambda \in \mathbf{C} - \mathbf{R}\) ; it is integrable on \(\mathbf{C}\) since \(\widetilde{f}\) has compact support. One deduces from Lebesgue's theorem (INT, IV, p. 137, § 3, n° 7, th. 6) and from the Helffer--Sjöstrand formula applied to \(f\) we deduce that \(\varpi(f(u_n))\) converges to \(\varpi(f(u))\) , hence \(p(f(u_n) - f(u))\) tends to 0.

Suppose now that \(f \in \mathcal{C}_0(\mathbf{R})\) Let \(\varepsilon > 0\) . There exists a function \(f_{\varepsilon}\) in \(\mathcal{D}(\mathbf{R})\) such that \(\| f_{\varepsilon} - f \|_{\infty} \leqslant \varepsilon\) (cf. prop. 4, a) of IV, p. 202 and INT, III, p. 45, § 1, n° 2, prop. 3). According to prop. 5, c) of IV, p. 275, one then has \(\| f(u) - f_{\varepsilon}(u) \| \leqslant \varepsilon\) and \(\| f(u_n) - f_{\varepsilon}(u_n) \| \leqslant \varepsilon\) for all \(n \in \mathbf{N}\) . Consequently, it follows that

\[
p (f (u _ {n}) - f (u)) \leqslant 2 c \varepsilon + p (f _ {\varepsilon} (u _ {n}) - f _ {\varepsilon} (u))
\]

for all \(n \in \mathbf{N}\) . Since \(f_{\varepsilon} \in \mathcal{D}(\mathbf{R})\) , the sequence \((p(f_{\varepsilon}(u_n) - f_{\varepsilon}(u)))_{n \in \mathbf{N}}\) converges to 0 by the preceding case, and therefore \(p(f(u_n) - f(u))\) converges to 0. This concludes the proof of a).

Let us prove assertion \(b\) ). Under the hypothesis concerning \(\mathcal{T}\) , it follows from TG, IX, p. 17, prop. 10 and the following remark, and from assertion \(a\) ), that the map \(u \mapsto f(u)\) of \(\mathcal{A}(\mathrm{E})\) in \(\mathcal{L}(\mathrm{E})\) is continuous when \(f \in \mathcal{C}_0(\mathbf{R})\) . Since the maps \(f \mapsto f(u)\) of \(\mathcal{C}_0(\mathbf{R})\) in \(\mathcal{L}(\mathrm{E})\) are continuous with norm \(\leqslant 1\) for all \(u \in \mathcal{A}(\mathrm{E})\) (prop. 5, \(c\) ) of IV, p. 275), assertion \(b\) ) follows from TG, X, p. 13, cor. 3.

Examples. --- 1) Let \(\mathscr{S}\) be a set of bounded subsets of E. The \(\mathscr{S}\)-topology on \(\mathcal{L}(\mathrm{E})\) (EVT, III, p. 13) is a locally convex topology less fine than the Banach space topology of \(\mathcal{L}(\mathrm{E})\) .

\begin{enumerate}
\def\labelenumi{\arabic{enumi})}
\setcounter{enumi}{1}
\tightlist
\item
  Let \(T_{b}\) the Banach space topology of \(\mathcal{L}(\mathrm{E})\) . For the topology \(T_{b}\) -resolvent, every \(u \in \mathcal{A}(\mathrm{E})\) admits a countable fundamental system of neighborhoods.
\end{enumerate}

Indeed, it suffices to show that for every \(\lambda \in \mathbf{C} - \mathbf{R}\) and every \(\varepsilon > 0\) , there exist \(\lambda' \in \mathbf{Q} + i\mathbf{Q}^*\) and an integer \(n \geqslant 1\) such that every partial operator \(v \in \mathcal{A}(\mathrm{E})\) satisfying \(\| \mathrm{R}(v, \lambda') - \mathrm{R}(u, \lambda')\| < 1/n\) also satisfies condition \(\|R(v,\lambda)-R(u,\lambda)\|<\varepsilon\) ; by writing

\[
\begin{array}{c} \| \mathrm{R} (v, \lambda) - \mathrm{R} (u, \lambda) \| \leqslant \| \mathrm{R} (v, \lambda) - \mathrm{R} (v, \lambda^ {\prime}) \| \\ + \| \mathrm{R} (v, \lambda^ {\prime}) - \mathrm{R} (u, \lambda^ {\prime}) \| + \| \mathrm{R} (u, \lambda^ {\prime}) - \mathrm{R} (u, \lambda) \|, \end{array}
\]

this property follows from formula (8) of IV, p. 245, from the estimates of proposition 17 of IV, p. 248, and from the fact that \(\mathbf{Q} + i\mathbf{Q}^*\) is everywhere dense in \(\mathbf{C} - \mathbf{R}\) .

The conclusion of the proposition is not valid in general if \(\mathcal{C}_{0}(\mathbf{R})\) is replaced by \(\mathcal{C}_{b}(\mathbf{R})\) (exercise 29 of IV, p. 366, e)). Nevertheless, one has the following result.

COROLLARY. --- Let \(\mathcal{T}_{s}\) the topology of simple convergence on \(\mathcal{L}(\mathrm{E})\) Let \(f \in \mathcal{C}_{b}(\mathbf{R})\) . For every sequence \((u_{n})_{n \in \mathbf{N}}\) in \(\mathcal{A}(\mathrm{E})\) which converges to \(u\) for the topology \(\mathcal{T}_{s}\) -resolvent, the sequence \((f(u_{n}))_{n \in \mathbf{N}}\) converges to \(f(u)\) in \(\mathcal{L}(\mathrm{E})\) endowed with the topology \(\mathcal{T}_{s}\) .

Let \((u_{n})\) be a sequence in \(\mathcal{A}(\mathrm{E})\) which converges to u for the topology \(T_{s}\) -resolvent. Let \(x \in E\) and let \(\varepsilon > 0\) .

For every integer \(N \in \mathbf{N}\) , let \(\varphi_{N} \in \mathcal{K}(\mathbf{R})\) be a function with support contained in \([-(N+1), N+1]\) such that \(0 \leqslant \varphi_{N} \leqslant 1\) and \(\varphi_{N}(t) = 1\) for all \(t \in [-N, N]\) . The functions \(\varphi_{N}\) converge pointwise to 1 as N tends to infinity, and satisfy \(|\varphi_{N}| \leqslant 1\) ; hence prop. 6 of IV, p. 276 implies that there exists \(N \in \mathbf{N}\) such that \(\|\varphi_{N}(u)x - x\| \leqslant \varepsilon\) . Set \(f_{N} = f\varphi_{N}\) .

For every \(n\in \mathbf{N}\) , we have

\[
\begin{array}{r l} \| f (u _ {n}) x - f (u) x \| \leqslant & \| f (u _ {n}) x - f _ {\mathrm{N}} (u _ {n}) x \| + \\ & \| f _ {\mathrm{N}} (u _ {n}) x - f _ {\mathrm{N}} (u) x \| + \| f _ {\mathrm{N}} (u) x - f (u) x \|. \end{array} \tag {15}
\]

For every \(v\in \mathcal{A}(\mathrm{E})\) , we have

\[
\begin{array}{r l} & {\| f (v) x - f _ {\mathrm{N}} (v) x \| = \| (f (1 - \varphi_ {\mathrm{N}})) (v) x \|} \\ & {\qquad \leqslant \| f (v) \|   \| x - \varphi_ {\mathrm{N}} (v) x \| \leqslant \| f \| _ {\infty}   \| x - \varphi_ {\mathrm{N}} (v) x \|} \end{array}
\]

(prop. 5, c) of IV, p. 275). Since \(\varphi_{\mathrm{N}} \in \mathcal{C}_0(\mathbf{R})\) , assertion \(a\) ) of prop. 11 applied to the topology \(\mathcal{T}_s\) -resolvent implies that \(\varphi_{\mathrm{N}}(u_n)x\) converges to \(\varphi_{\mathrm{N}}(u)x\) as \(n \to +\infty\) . Consequently, for every \(n\) sufficiently large, one has

\[
\left\| x - \varphi_ {\mathrm{N}} (u _ {n}) x \right\| \leqslant \left\| x - \varphi_ {\mathrm{N}} (u) x \right\| + \varepsilon \| x \| \leqslant (1 + \| x \|) \varepsilon .
\]

Inequality (15) then becomes

\[
\| f (u _ {n}) x - f (u) x \| \leqslant \| f \| _ {\infty} (2 + \| x \|) \varepsilon + \| f _ {\mathrm{N}} (u _ {n}) x - f _ {\mathrm{N}} (u) x \|
\]

for all sufficiently large n. One concludes the proof by means of assertion a) of loc. cit., applied to the function \(f_{\mathrm{N}} \in \mathcal{C}_{0}(\mathbf{R})\) and to the topology \(T_{s}\) -resolvent.

\subsection{Polar decomposition}

Lemma 8. --- Let E be a separated topological vector space. Let \((\mathrm{E}_{1},\|\cdot\|_{1})\) and \((\mathrm{E}_{2},\|\cdot\|_{2})\) be dense subspaces of E endowed with Hilbert structures such that the canonical injections of \(E_{1}\) and \(E_{2}\) into E are continuous. Let F be a subspace of \(E_{1}\cap E_{2}\) , dense in \(E_{1}\) and \(E_{2}\) for these Hilbert structures. If \(\|x\|_{1}=\|x\|_{2}\) for all \(x\in F\) , then \(E_{1}=E_{2}\) and \(\|\cdot\|_{1}=\|\cdot\|_{2}\) .

Let \(x \in E_{1}\) . There exists a sequence \((x_{n})_{n \in \mathbf{N}}\) in F such that \(x_{n}\) converges to x in the Hilbert space \(E_{1}\) . Since \(\|x_{n} - x_{m}\|_{2} = \|x_{n} - x_{m}\|_{1}\) for all integers n and m, the sequence \((x_{n})\) is a Cauchy sequence in \(E_{2}\) . Let y be its limit. We have \(\|y\|_{2} = \lim\|x_{n}\|_{2} = \|x\|_{1}\) . Since the canonical injections of \(E_{1}\) and \(E_{2}\) into E are continuous, we have \(x_{n} \to x\) in E and likewise \(x_{n} \to y\) in E. Thus, it follows that x = y and \(\|x\|_{1} = \|x\|_{2}\) . In particular, \(E_{1} \subset E_{2}\) . We conclude by symmetry.

Let E be a complex Hilbert space. Let u be a positive self-adjoint partial operator on E. For every \(\alpha \in \mathbf R_{+}\) we set \(u^{\alpha} = f(u)\) , where f is the continuous map \(x \mapsto x^{\alpha}\) of \(\mathbf R_{+}\) into \(\mathbf R\). This is a positive self-adjoint partial operator (cor. 1, a) of IV, p. 273). When \(u \in \mathcal{L}(E)\) , this notation is compatible with the notation relative to the C*-algebra \(\mathcal{L}(E)\) (cf. prop. 16 of I, p. 118). If \(\alpha\) is a positive integer, the partial operator \(u^{\alpha}\) coincides with the partial operator defined by composition \(u \circ \cdots \circ u\) (cor. 1, b) of IV, p. 273).

Let \(\beta \in \mathbf{R}_{+}\) . We have \(u^{\alpha \beta} = (u^{\alpha})^{\beta}\) (cor. 4 of IV, p. 274). In particular, if \(\alpha > 0\) , then the partial operator \(u^{1 / \alpha}\) is the unique positive self-adjoint partial operator \(v\) on E such that \(v^{\alpha} = u\) .

Suppose that \(0 \leqslant \alpha \leqslant \beta\) . We then have \(\operatorname{dom}(u^{\beta}) \subset \operatorname{dom}(u^{\alpha})\) : indeed, for every real number \(x \geqslant 0\) , we have \(x^{\alpha} \leqslant 1 + x^{\beta}\) , and the result then follows from the definition of the domain of \(u^{\alpha}\) (cf. prop. 3 of IV, p. 271).

Let u be a closed densely defined operator from E into a complex Hilbert space F. The partial operator \(u^{*}u\) is self-adjoint and positive (prop. 12 of IV, p. 241). We denote \(|u| = (u^{*}u)^{1/2}\) .

PROPOSITION 12. --- Let u be a closed densely defined operator from E into a complex Hilbert space F.

\begin{enumerate}
\def\labelenumi{\alph{enumi})}
\item
  The domain of the positive self-adjoint partial operator \(|u|\) coincides with the domain of \(u\) ;
\item
  There exists a unique partially isometric linear map \(j\) from E into F such that \(u = j|u|\) and \(\operatorname{Ker}(j) = \operatorname{Ker}(u)\) ;
\item
  Let \(u_{1}\) be a positive self-adjoint operator on E and \(j_{1}\) a partially isometric linear map from E into F such that \(u = j_{1}u_{1}\) and \(\operatorname{Ker}(j_{1}) = \operatorname{Ker}(u_{1})\) . We then have \(u_{1} = |u|\) and \(j_{1} = j\) .
\end{enumerate}

The domain of \(u^{*}u\) is contained in \(\operatorname{dom}(u)\) and in \(\operatorname{dom}(|u|)\) . It is dense in the Hilbert spaces \(E_{u}\) (loc. cit.) and \(E_{|u|}\) (cor. 1, e) of IV, p. 273) and, for every \(x \in \operatorname{dom}(u^{*}u)\) , we have

\[
\begin{array}{c} \| x \| _ {u} ^ {2} = \| x \| ^ {2} + \langle u (x) | u (x) \rangle = \| x \| ^ {2} + \langle x | (u ^ {*} u) (x) \rangle \\ = \| x \| ^ {2} + \langle | u | (x) | | u | (x) \rangle = \| x \| _ {| u |} ^ {2}. \end{array}\tag{16}
\]

Consequently, the Hilbert spaces \(E_{u}\) and \(E_{|u|}\) are equal (lemma 8), hence \(\text{dom}(u) = \text{dom}(|u|)\) and \(\|u(x)\| = \||u|(x)\|\) for all \(x \in \text{dom}(u)\) .

Formula (16) implies that \(\operatorname{Ker}(u) = \operatorname{Ker}(|u|)\) and that there exists a unique isometric linear map \(v\) from \(\overline{\operatorname{Im}(|u|)}\) onto \(\overline{\operatorname{Im}(u)}\) satisfying \(v(|u|(x)) = u(x)\) for all \(x \in \operatorname{dom}(|u|)\) . Since \(|u|\) is self-adjoint, we have \(\operatorname{Im}(|u|)^{\circ} = \operatorname{Ker}(|u|)\) (prop. 7, c) of IV, p. 236). There therefore exists a unique partial isometry \(j\) from E into F which extends \(v\) and vanishes on \(\operatorname{Ker}(|u|) = \operatorname{Ker}(u)\) . Since \(\operatorname{E} = \operatorname{Ker}(u) \oplus \overline{\operatorname{Im}(|u|)}\) , this map is the unique partially isometric map such that \(u = j|u|\) and \(\operatorname{Ker}(j) = \operatorname{Ker}(u)\) .

Let us prove c). We have \(u_{1}j_{1}^{*}j_{1}u_{1}\subset(j_{1}u_{1})^{*}j_{1}u_{1}=u^{*}u\) . The linear map \(j_{1}^{*}j_{1}\) is the orthogonal projector with kernel \(\mathrm{Ker}(j_{1})=\mathrm{Ker}(u_{1})\) (EVT, V, p. 41, prop. 5 (ii)) and therefore with image \(\mathrm{Ker}(u_{1})^{\circ}=\overline{\mathrm{Im}(u_{1})}\) (prop. 7, c) of IV, p. 236). Consequently, \(u_{1}^{2}\subset u^{*}u\) , whence \(u_{1}^{2}=u^{*}u\) since these two operators are self-adjoint. Thus, it follows that \(u_{1}=(u^{*}u)^{1/2}\) and the uniqueness assertion of b) ultimately shows that \(j_{1}=j\) .

DEFINITION 7. --- Let u be a closed operator with dense domain from E into a complex Hilbert space F. The pair \((j, |u|)\) determined by prop. 12 is called the polar decomposition of u.

Suppose that \(u \in \mathcal{L}(\mathrm{E};\mathrm{F})\) . Its polar decomposition in the sense of this definition coincides with that of Definition 4 in I, p. 140.

\subsection{Self-adjoint operators defined by a positive partial Hermitian form}

In this issue, E denotes a complex Hilbert space.

DEFINITION 8. --- Let D be a vector subspace of E. A partial Hermitian form on E with domain D is a correspondence \(q = (\Gamma, \mathrm{E} \times \mathrm{E}, \mathbf{C})\) whose domain of definition is \(\mathrm{D} \times \mathrm{D}\) , whose graph \(\Gamma\) is functional, and such that the map from \(\mathrm{D} \times \mathrm{D}\) into \(\mathbf{C}\) defined by \(\Gamma\) is a Hermitian form. We denote \(\operatorname{dom}(q) = \mathrm{D}\) .

We say that \(q\) is a positive partial form if the Hermitian form it defines is positive.

Let u be a partial self-adjoint operator on E. For all elements x and y of E, we denote \(\mu_{x,y}\) (resp. \(\mu_{x}\) ) as the spectral measure of \((x,y)\) (resp. of x) relative to u.

Let \((j, |u|)\) be the polar decomposition of \(u\) (definition 7 of IV, p. 290). Let D' denote the domain of \(|u|^{1/2}\) . By definition, it is the space of \(x \in E\) such that the function \(z \mapsto |z|\) is integrable on \(\mathrm{Sp}(u)\) with respect to the measure \(\mu_x\) . It contains \(\mathrm{dom}(u)\) .

Let \((x,y)\in\mathrm{D}^{\prime}\times\mathrm{D}^{\prime}\) . The identity function of \(\mathrm{Sp}(u)\) is integrable with respect to the measure \(\mu_{x,y}\) (prop. 4, b) of IV, p. 271); we set

\[
q _ {u} (x, y) = \int_ {\operatorname{Sp} (u)} t d \mu_ {x, y} (t).
\]

If \(y \in \operatorname{dom}(u)\) , we have \(q_u(x, y) = \langle x \mid u(y) \rangle\) by formula (6) of IV, p. 271. The map \(q_u\) is a Hermitian form on D'. It defines a partial Hermitian form on E, said to be associated with u. If the operator u is positive, then the form \(q_u\) is a positive partial Hermitian form.

DEFINITION 9. --- Let \(q\) be a positive partial Hermitian form on E. We denote by \(\mathrm{E}_q\) the separated pre-Hilbert space \(\operatorname{dom}(q)\) equipped with the scalar product

\[
(x \mid y) _ {q} = \langle x \mid y \rangle + q (x, y).
\]

We denote by \(\|x\|_{q}\) the norm of \(x \in E_{q}\) . The form q is said to be closed if \(E_{q}\) is a Hilbert space.

PROPOSITION 13. --- Let u be a positive self-adjoint partial operator on E. Let \(q_{u}\) be the positive partial form associated with u.

\begin{enumerate}
\def\labelenumi{\alph{enumi})}
\item
  The domain of \(q_{u}\) is the domain of \(u^{1/2}\) , and for all x and y in \(\operatorname{dom}(u^{1/2})\) , we have \(q_{u}(x,y)=\langle u^{1/2}(x)\mid u^{1/2}(y)\rangle\) ;
\item
  The positive partial form \(q_{u}\) is closed. The domain of u is a core for \(q_{u}\) .
\end{enumerate}

Since \(u\) is positive, we have \(|u| = u\) . The domain \(\mathrm{dom}(|u|^{1/2})\) of \(q_u\) therefore coincides with \(\mathrm{dom}(u^{1/2})\) and we have

\[
\begin{array}{c} q _ {u} (x, y) = \int_ {\operatorname{Sp} (u)} t d \mu_ {x, y} (t) = \int_ {\operatorname{Sp} (u)} t ^ {1 / 2} \cdot t ^ {1 / 2} d \mu_ {x, y} (t) \\ = \langle u ^ {1 / 2} (x) | u ^ {1 / 2} (y) \rangle \end{array}
\]

for \(x\) and \(y\) in \(\mathrm{dom}(q_u)\) (prop. 4, b) of IV, p. 271). This proves the assertion \(a\) ).

Moreover, the pre-Hilbert space \(E_{q_{u}}\) then coincides with the pre-Hilbert space \(E_{u^{1/2}}\) associated with \(u^{1/2}\) (def. 4 of IV, p. 230). Since \(u^{1/2}\) is a closed partial operator, this space is a Hilbert space (prop. 4 of IV, p. 230) and \(\text{dom}(u)\) is dense in \(E_{u^{1/2}}\) by assertion \(e\) ) of Cor. 1 of IV, p. 273.

Let \(q\) be a partial Hermitian form on \(E\) with domain \(D\) dense in \(E\) . For \(y \in D\) , denote by \(\lambda_y\) the linear form \(x \mapsto q(y, x)\) on \(D\) . We denote by \(\widetilde{D}\) the set of \(y \in D\) such that the linear form \(\lambda_y\) is continuous on \(D\) .

Let \(y \in \widetilde{D}\) . Since D is dense in E, there exists a unique continuous linear form on E extending \(\lambda_{y}\) . We still denote it by \(\lambda_{y}\) . There exists a unique element \(u(y)\) in E such that \(\lambda_{y}(x) = \langle u(y) | x \rangle\) for all \(x \in E\) (EVT, V, p. 15, th. 3). The map \(y \mapsto u(y)\) is linear from \(\widetilde{D}\) into E; we say that the partial operator u on E with domain \(\widetilde{D}\) is the partial operator representing q. We therefore have

\[
q (x, y) = \langle x \mid u (y) \rangle
\]

for \(y \in \operatorname{dom}(u)\) and \(x \in D\) .

Remark. --- Let \(q\) be a closed positive partial form and \(q'\) a continuous positive Hermitian form on E. The positive Hermitian form defined on \(\text{dom}(q)\) by \((x, y) \mapsto q(x, y) + q'(x, y)\) is a closed positive partial Hermitian form with the same domain as \(q\) . We denote it by \(q + q'\) .

According to EVT, V, p. 16, cor. 2, there exists a unique linear map \(u' \in \mathcal{L}(\mathrm{E})\) such that \(q'(x,y) = \langle x \mid u'(y) \rangle\) for all \((x,y) \in \mathrm{E} \times \mathrm{E}\) . The endomorphism \(u'\) is positive and the partial operator representing the closed positive partial Hermitian form \(q + q'\) is \(u + u'\) .

PROPOSITION 14. --- Let q be a closed positive partial form with dense domain on E and let u be the partial operator representing q.

\begin{enumerate}
\def\labelenumi{\alph{enumi})}
\item
  The domain of u is dense in the Hilbert space \(E_{q}\) ;
\item
  The partial operator u is self-adjoint and positive.
\end{enumerate}

Since q is closed, the pre-Hilbert space \(E_{q}\) of Definition 9 is a Hilbert space.

Let us show that the partial operator \(u + 1_{E}\) with domain \(\operatorname{dom}(u)\) is bijective. Since

\[
\langle x \mid (u + 1 _ {\mathrm{E}}) (x) \rangle = q (x, x) + \| x \| ^ {2} \geqslant \| x \| ^ {2}
\]

for all \(x \in \operatorname{dom}(u)\) , this partial operator is injective.

Let \(y \in E\) . The linear form on \(E_{q}\) defined by \(x \mapsto \langle y | x \rangle\) is continuous since \(\|x\| \leqslant \|x\|_{q}\) . There therefore exists \(y' \in E_{q}\) such that

\[
\langle y \mid x \rangle = (y ^ {\prime} \mid x) _ {q} = \langle y ^ {\prime} \mid x \rangle + q (y ^ {\prime}, x)
\]

for all \(x \in E_{q}\) (EVT, V, p. 15, th. 3). By definition, this means that \(y'\) belongs to the domain of the partial operator \(\widetilde{u}\) representing the partial form \((x, y) \mapsto (x \mid y)_{q}\) with domain dom(q), and that \(\widetilde{u}(y') = y\) . Since \(\widetilde{u} = u + 1_{E}\) by the remark above, we deduce from this that the partial operator \(u + 1_{E}\) is also surjective, hence bijective.

Let us prove that the domain of \(u\) is dense in \(\mathrm{E}_q\) . Let \(y \in \mathrm{E}_q\) be orthogonal to \(\operatorname{dom}(u)\) in \(\mathrm{E}_q\) . There exists \(y' \in \operatorname{dom}(u)\) such that \(u(y') + y' = y\) . We then have

\[
0 = (y \mid y ^ {\prime}) _ {q} = \langle y \mid y ^ {\prime} \rangle + q (y, y ^ {\prime}) = \langle y \mid y ^ {\prime} + u (y ^ {\prime}) \rangle = \| y \| ^ {2}
\]

hence \(y = 0\) . The assertion \(a\) ) is therefore proved.

Since \(\mathrm{dom}(q)\) is dense in E and \(\| x\| \leqslant \| x\| _q\) for all \(x\in \mathrm{dom}(q)\) , assertion \(a\) ) implies that \(\mathrm{dom}(u)\) is dense in E.

Since the form \(q\) is Hermitian (resp. positive), for all \(x\) and \(y\) in \(\operatorname{dom}(u)\) , we have

\[
\langle y \mid u (x) \rangle = q (y, x) = \overline {{q (x , y)}} = \overline {{\langle x \mid u (y) \rangle}} = \langle u (y) \mid x \rangle ,
\]

(resp. \(\langle x \mid u(x) \rangle = q(x, x) \geqslant 0\) ) so that u is symmetric (resp. positive). Finally, the partial operator \(u + 1_{E}\) is self-adjoint by the corollary of Proposition 10 of IV, p. 240, and the same is true of u.

THEOREM 3. --- The map that associates the positive partial form \(q_u\) with a positive self-adjoint operator u on E is a bijection between the set of positive self-adjoint partial operators on E and the set of closed positive partial forms with dense domain on E. The inverse bijection associates with a positive partial form q the partial operator representing q.

By Propositions 13, b) and 14, b), the maps described in the statement are well defined. Let us show that they are inverse bijections of one another.

Let \(u\) be a positive self-adjoint partial operator on E and \(q\) the positive partial form associated with \(u\) . Denote by \(v\) the positive self-adjoint operator representing \(q\) . Let \(y \in \mathrm{dom}(u) \subset \mathrm{dom}(u^{1/2}) = \mathrm{dom}(q)\) . For every \(x \in \mathrm{dom}(q) = \mathrm{dom}(u^{1/2})\) , we have

\[
q (y, x) = \langle u ^ {1 / 2} (y) | u ^ {1 / 2} (x) \rangle = \langle u (y) | x \rangle ,
\]

hence y belongs to the domain of v and satisfies \(v(y) = u(y)\) . The partial operator v is therefore an extension of u; since u and v are self-adjoint, they are equal.

Conversely, let q be a closed positive partial form with dense domain and let u be the positive self-adjoint partial operator representing q. We have \(\operatorname{dom}(u) \subset \operatorname{dom}(u^{1/2})\) . For x and y in \(\operatorname{dom}(u)\) , we have

\[
q _ {u} (x, y) = \langle u ^ {1 / 2} (x) | u ^ {1 / 2} (y) \rangle = \langle x | u (y) \rangle = q (x, y)
\]

by Proposition 13, a). Thus, the Hilbert spaces \(\mathrm{E}_q\) and \(\mathrm{E}_{q_u}\) both contain \(\operatorname{dom}(u)\) as a dense subspace (Prop. 14, a) and Prop. 13, b), respectively) and \(\|x\|_q = \|x\|_{q_u}\) for all \(x \in \operatorname{dom}(u)\) . It follows that \(\mathrm{E}_q = \mathrm{E}_{q_u}\) and that \(q = q_u\) (lemma 8).

COROLLARY. --- Let \(q\) be a closed positive partial form on \(E\) and \(u\) the positive self-adjoint operator representing \(q\) . The domain of \(q\) is equal to the domain of \((1_{E} + u)^{1/2}\) , and we have

\[
\| x \| _ {q} = \| (1 _ {\mathrm{E}} + u) ^ {1 / 2} x \|
\]

for all \(x \in \operatorname{dom}(q)\) .

The domain of \(q\) is equal to the domain of \(u^{1/2}\) (prop. 13, \(a\) ), which coincides with that of \((1_{\mathrm{E}} + u)^{1/2}\) (cf. prop. 3 of IV, p. 271). For every \(x \in \operatorname{dom}(u) \subset \operatorname{dom}(u^{1/2})\) , we have

\[
\| (1 _ {\mathrm{E}} + u) ^ {1 / 2} x \| ^ {2} = \langle x | (1 _ {\mathrm{E}} + u) (x) \rangle = \| x \| ^ {2} + \langle x | u (x) \rangle = \| x \| _ {q} ^ {2}.
\]

As the domain of \(u\) is dense in the Hilbert space \(\mathrm{E}_q\) (prop. 14, a)), this formula extends by continuity to all \(x \in \mathrm{dom}(u^{1/2})\) .

Example. --- Let u be a positive partial operator on E which is not necessarily closed. One defines a positive partial form q with domain dom(u) by

\[
q (x, y) = \langle x \mid u (y) \rangle
\]

for \(x\) and \(y\) in \(\mathrm{dom}(u)\) . Let \(\mathrm{E}_q\) be the separated pre-Hilbert space of Definition 9 of IV, p. 292. Let \(\widetilde{\mathrm{E}}_q\) be the completed Hilbert space of \(\mathrm{E}_q\) , still denoted \((x,y) \mapsto (x \mid y)_q\) for the scalar product. Since the canonical inclusion \(\iota\) of \(\mathrm{E}_q\) in E is continuous, it admits a unique continuous extension, denoted \(\widetilde{\iota}\) , of \(\widetilde{\mathrm{E}}_q\) in E. The Hermitian form \(q\) is continuous on \(\mathrm{E}_q\) , hence also extends to a unique continuous positive Hermitian form \(\widetilde{q}\) on \(\widetilde{\mathrm{E}}_q\) .

Let us prove that the linear map \(\widetilde{\iota}\) is injective. Let \(x\in\mathrm{Ker}(\widetilde{\iota})\) . Consider a sequence \((x_{n})_{n\in\mathbf{N}}\) in \(\mathrm{E}_{q}\) which converges to \(x\) in \(\widetilde{\mathrm{E}}_{q}\) . We then have

\[
\lim _ {n \to + \infty} \iota (x _ {n}) = \widetilde {\iota} (x) = 0,
\]

therefore the sequence \((x_{n})_{n\in\mathbf{N}}\) converges to 0 in E. Let \(y\in E_{q}\) . Since \(\widetilde{q}\) is continuous on \(\widetilde{E}_{q}\) , it follows that

\[
\begin{array}{r l} (x \mid y) _ {q} = \lim _ {n \to + \infty} (x _ {n} \mid y) _ {q} = & \lim _ {n \to + \infty} \Big (\langle x _ {n} \mid y \rangle + q (x _ {n}, y) \Big) \\ & = \lim _ {n \to + \infty} \Big (\langle x _ {n} \mid y \rangle + \langle x _ {n} \mid u (y) \rangle \Big) = 0. \end{array}
\]

As the space \(E_{q}\) is dense in \(\widetilde{E}_{q}\) , we deduce that x = 0, as desired.

Identifying \(\widetilde{E}_{q}\) with its image under \(\widetilde{\iota}\) in E, we interpret \(\widetilde{q}\) as a closed positive partial form with dense domain extending q. The positive self-adjoint operator associated with \(\widetilde{q}\) is a self-adjoint extension of u. It is called the Friedrichs extension of u.

Note. --- Let \(c \in \mathbf R_{+}\) . Let q be a closed, densely defined partial Hermitian form such that \(q(x, x) \geqslant -c \|x\|^2\) for all x belonging to the domain of q. This means that the closed partial Hermitian form with domain dom(q) defined by \(\widetilde{q}(x, y) = q(x, y) + c \langle x \mid y \rangle\) is a positive partial form. By Theorem 3, such forms therefore correspond to self-adjoint partial operators u on E such that \(u + c1_{E}\) is positive, that is, such that u is bounded below (def. 7 of IV, p. 237) by -c.

Conversely, let u be a symmetric partial operator on E such that

\[
\langle x \mid u (x) \rangle \geqslant - c \| x \| ^ {2}\tag{17}
\]

for all \(x \in \operatorname{dom}(u)\) . Then \(v = u + c1_{E}\) is a positive partial operator on E. The Friedrichs extension of u is the self-adjoint operator \(\widetilde{v} - c1_{E}\) , where \(\widetilde{v}\) is the Friedrichs extension of v; it is a self-adjoint extension of u that does not depend on the choice of the real number c satisfying (17).